\section{Affine-invariant geometry: classification details}
\label{app:affine-classification}

This appendix records the representation argument behind the metric
classification in \Cref{sec:geometry}.  The affine group acts transitively on
\(\R^d\times\SPD(d)\).  It is therefore enough to classify inner products on
the tangent space at \(o=(0,\Id)\) that are invariant under the isotropy group
\(O(d)\).

At \(o\), the tangent representation is
\[
  \R^d\oplus\Sym(d)
  =
  \R^d\oplus\operatorname{span}\{\Id\}\oplus\Sym_0(d),
\]
where \(\Sym_0(d)\) denotes the traceless symmetric matrices.  These three
summands are pairwise inequivalent irreducible \(O(d)\)-representations
(with the familiar low-dimensional interpretation when \(d=1\)).
Schur's lemma therefore leaves one scalar weight on each summand.
Equivalently, every invariant quadratic form is
\[
  \eta\norm u^2+\omega\Tr(X^2)+\tau(\Tr X)^2.
\]
There is no mean--covariance term.  One can see this without representation
terminology: the isotropy element \(-\Id\) sends \(u\) to \(-u\) and fixes
\(X\), so any cross term changes sign and must vanish.

Transporting this form from \(o\) to \(a=(m,C)\) by
\((u,X)\mapsto(C^{-1/2}u,C^{-1/2}XC^{-1/2})\) gives
\[
  g_a((u,X),(v,Y))
  =
  \eta u^\top C^{-1}v
  +\omega\Tr(C^{-1}XC^{-1}Y)
  +\tau\Tr(C^{-1}X)\Tr(C^{-1}Y).
\]
Positivity on the mean and traceless summands requires
\(\eta>0\) and \(\omega>0\).  On \(X=s\Id\), the covariance part is
\(d(\omega+d\tau)s^2\), hence
\(\tau>-\omega/d\).  These conditions are also sufficient.

For completeness, let
\[
  g:=-\E_{\N(m,C)}\nabla V,\qquad
  A:=\E_{\N(m,C)}\nabla^2V,\qquad
  R:=C^{-1}-A.
\]
The Euclidean differentials of the objective are
\[
  D_m\calE[u]=-g^\top u,\qquad
  D_C\calE[X]=-\frac12\Tr(RX).
\]
Solving
\(g_a(\grad\calE,(u,X))=D\calE[(u,X)]\) separately on the mean,
traceless, and trace covariance components gives
\begin{align*}
  \dot m
  &=\frac1\eta Cg,\\
  \dot C
  &=
  \frac1{2\omega}CRC
  -
  \frac{\tau}{2\omega(\omega+d\tau)}
  \Tr(CR)\,C.
\end{align*}
Whitening a Gaussian target at its optimizer yields
\(\dot m=-m/\eta+O(\norm{a-a_\star}^2)\).
For \(C=\Id+X\), the traceless covariance part obeys
\(\dot X_0=-X_0/(2\omega)+O(\norm X^2)\), while the trace part obeys
\[
  \frac{\dd}{\dd t}\frac{\Tr X}{d}
  =
  -\frac1{2(\omega+d\tau)}
  \frac{\Tr X}{d}
  +O(\norm X^2).
\]
Balancing all three linear rates, modulo an overall multiplication of the
metric, requires
\[
  \eta=1,\qquad \omega=\frac12,\qquad \tau=0,
\]
which is the Fisher--Rao metric.
